\documentclass[10pt,a4paper]{amsart}
\usepackage[margin=19mm]{geometry}
\usepackage{amsmath,amssymb,mathtools}
\usepackage[scr=rsfs,cal=euler]{mathalfa}
\usepackage{xfrac}
\usepackage[unicode,hidelinks,bookmarks=false]{hyperref}
\newcommand{\ie}{i.e.\ }
\newcommand{\cf}{cf.\ }
\newcommand{\resp}{respectively\ }
\newcommand{\Subsection}{\subsection}
\providecommand{\nobreakdash}{\nobreak\hbox{-}}
\setlength{\emergencystretch}{2em}
\multlinegap0pt
\usepackage{amssymb}
\usepackage{mathtools}
\usepackage[shortlabels]{enumitem}
\usepackage{xcolor}
\newtheorem{lemma}{Lemma}
\usepackage{cleveref}
\crefname{lemma}{Lemma}{Lemmas}
\newcommand{\abs}[1]{\left\lvert #1\right\rvert}
\newcommand{\cA}{\mathcal A}
\newcommand{\dd}{\,d}
\newcommand{\norm}[1]{\left\lVert #1\right\rVert}
\title{Common Causal It\^o Flows under Nondominated Martingale Laws: Capacity Cores and Robust Sensitivities}
\author{Guangqian Zhao}
\date{Source: arXiv:2509.13675v8 (CC BY 4.0)}
\begin{document}
\maketitle

\noindent\emph{Source and licence.} Common Causal It\^o Flows under Nondominated Martingale Laws: Capacity Cores and Robust Sensitivities. Version arXiv:2509.13675v8 (CC BY 4.0), distributed by arXiv under the Creative Commons Attribution 4.0 International licence, http://creativecommons.org/licenses/by/4.0/, as recorded in the arXiv licence metadata for this exact version and shown on the abstract page https://arxiv.org/abs/2509.13675v8. Full licence notice: \texttt{licenses/arxiv-2509.13675-CC-BY-4.0.txt}.\par\smallskip

\noindent\textit{Editorial scope.} Counted unit: the article's lemma lem:dyadic-estimate (source0.tex lines 660--671). The article's complete original proof of it is delivered with it (source0.tex lines 673--735, one \texttt{proof} environment of 63 lines). No mathematics is written, repaired or re-derived: the statement and the proof are the article's own lines, in the article's own notation, and no retained line is substituted or re-flowed.  Retained uncounted as the source-local context the proof needs: lines 531--535; lines 595--603; lines 654--658.  Nothing was omitted from any retained span, so the package carries no omission record.  No retained line contains a citation, so the wrapper carries no bibliography and no bibliography entry.  No retained line contains a figure environment, an image-inclusion command or drawing code, so no image is delivered and the package declares no asset dependency.  Editorial formatting only, in addition: an AMS \texttt{amsart} preamble that restates the article's own declarations and macros used by the retained lines, namely the environments \texttt{lemma} on separate counters, together with the standard packages those lines need.  The retained environments print in this document as Lemma 1 in order of appearance, whereas the article prints the same items with the numbering of its own full document.  The complete native source of the article as distributed in the arXiv e-print for this exact version is bundled unchanged as \texttt{sources/185188/source0.tex}.
\begin{equation}\label{eq:martingale-moment}
  \sup_{P\in\mathfrak M_\Lambda}
  E^P\bigl[\abs{X_t-X_s}^r\bigr]
  \le C_r\Lambda^{r/2}\abs{t-s}^{r/2}.
\end{equation}

\begin{lemma}[Causal Lipschitz envelope]\label{lem:lipschitz-envelope}
The map \(\mathcal R_\Lambda\) takes nonnegative increasing functions into
\(\cA_\Lambda\), is causal, and satisfies
\[
  \norm{\mathcal R_\Lambda a-\mathcal R_\Lambda b}_\infty
  \le\norm{a-b}_\infty.
\]
If \(a\in\cA_\Lambda\), then \(\mathcal R_\Lambda a=a\).
\end{lemma}

\begin{equation}\label{eq:qv-martingale-error}
  \widetilde V_t^n(X)-\langle X\rangle_t^P
  =
  2\int_0^t(X_r-X_{\underline r_n})\dd X_r.
\end{equation}

\begin{lemma}[Uniform dyadic estimate]\label{lem:dyadic-estimate}
For every \(r\ge2\), there is \(C_{r,\Lambda,T}<\infty\) such that
\begin{equation}\label{eq:qv-Lr-rate}
  \sup_{P\in\mathfrak M_\Lambda}
  \left(
    E^P\left[
      \norm{q^n(X)-\langle X\rangle^P}_\infty^r
    \right]
  \right)^{1/r}
  \le C_{r,\Lambda,T}2^{-n/2}.
\end{equation}
\end{lemma}

\begin{proof}
Let \(h_n=T2^{-n}\).  By BDG, \eqref{eq:qv-martingale-error}, and
\(\dd\langle X\rangle_r^P\le\Lambda\dd r\),
\[
  \left\|
    \norm{\widetilde V^n-\langle X\rangle^P}_\infty
  \right\|_{L^r(P)}
  \le
  C_r
  \left(
    E^P\left[
      \left(
        \int_0^T
        \abs{X_r-X_{\underline r_n}}^2
        \dd\langle X\rangle_r^P
      \right)^{r/2}
    \right]
  \right)^{1/r}.
\]
Jensen's inequality and Fubini's theorem give
\begin{align*}
E^P\left[
  \left(\int_0^T
    \abs{X_r-X_{\underline r_n}}^2\dd\langle X\rangle_r^P
  \right)^{r/2}
\right]
&\le
\Lambda^{r/2}T^{r/2-1}
\int_0^T E^P\!\left[
  \abs{X_r-X_{\underline r_n}}^r
\right]\dd r\\
&\le C_{r,\Lambda,T}h_n^{r/2},
\end{align*}
where \eqref{eq:martingale-moment} is used in the last step.
Consequently,
\[
  \left\|
    \norm{\widetilde V^n-\langle X\rangle^P}_\infty
  \right\|_{L^r(P)}
  \le C_{r,\Lambda,T}h_n^{1/2}.
\]

The difference between \(\widetilde V^n\) and the completed sum \(V^n\) is
bounded by the largest squared oscillation on one dyadic cell.  Hence
\begin{align*}
  E^P\left[
    \norm{\widetilde V^n-V^n}_\infty^r
  \right]
  &\le
  \sum_{k<2^n}
  E^P\left[
    \sup_{t_k^n\le t\le t_{k+1}^n}
    \abs{X_t-X_{t_k^n}}^{2r}
  \right]\\
  &\le C_{r,\Lambda,T}2^nh_n^r
   \le C_{r,\Lambda,T}h_n^{r-1},
\end{align*}
where the maximal BDG inequality is applied on each dyadic cell.  For
\(r\ge2\), the resulting \(L^r\)-norm is bounded by a constant times
\(h_n^{1/2}\).  Finally, \Cref{lem:lipschitz-envelope} is nonexpansive and
fixes \(\langle X\rangle^P\).
This proves \eqref{eq:qv-Lr-rate}.
\end{proof}

\end{document}
